%قبل از شروع مسئله ترجیحاً چیزی ننویسید. اما اگر خواستید بنویسید، در قالب لیستی (فقرات و شمارش) بنویسید.
%برای اضافه کردن نمره، بعدی از شروع‌مسئله، [x] را اضافه کنید و به جای x  نمره مورد نظر را قرار دهید. برای زیرمسئله‌ها هم به همین صورت است.
%\شروع‌مسئله[10]
\شروع‌مسئله [۱۲]
% صورت اصلی مسئله را در این‌جا بنویسید.
توابع زیر را به صورت مجانبی تحلیل کنید و به صورت کامل توضیح دهید. برای آشنایی بیشتر و مرور روش‌های تحلیل می‌توانید به بخش ۳-۵ کتاب داده‌ساختارها و مبانی الگوریتم‌ها مراجعه کنید.
\شروع{زیرمسائل}
\زیرمسئله $T(n) = T(\floor{\frac{n}{2}}) + T(\ceil{\frac{n}{2}}) + n$

\پاسخ{
		از حدس و استقرا برای حل این مسئله استفاده می‌کنیم.  $T \in \Theta(n\log(n))$ می‌تواند با توجه به رابطۀ بازگشتی، حدسی منطقی باشد. فرض می‌کنیم برای $b$ و $c$ و $n$ های کوچکتر از $m$ داریم:
	$$T(n) \leq bn\log(n) + c$$
	می‌خواهیم داشته باشیم:
	$$T(m) = T(\floor{\frac{m}{2}}) + T(\ceil{\frac{m}{2}}) + m \leq b(\floor{\frac{m}{2}}\log(\floor{\frac{m}{2}}) + \ceil{\frac{m}{2}}\log(\ceil{\frac{m}{2}})) + 2c + m \leq 
	bm\log(m) + c$$
	$$\Leftrightarrow b(\floor{\frac{m}{2}} +\ceil{\frac{m}{2}})\log(\frac{m+1}{2}) + c + m \leq bm\log(m) = bm\log(\frac{m+1}{2}) + c + m \leq bm\log(m)$$
	$$\Leftrightarrow 1 + \frac{c}{m} \leq b\log(\frac{2m}{m+1})$$
	برای این که نابرابری بالا برقرار باشد کافی است $c$ را منفی در نظر بگیریم و $b$ را به اندازه‌ی کافی بزرگ. در این صورت طبق استقرا برای هر $n$ خواهیم داشت:
	$$T(n) \leq bn\log(n) + c \Rightarrow T \in O(n\log(n))$$
	به همین صورت می‌توان ثابت کرد که $T \in \Omega(n\log(n))$. در نتیجه $T \in \Theta(n\log(n))$.
	
}
\زیرمسئله $T(n) = T(n - 1) + \frac{1}{n}$

\پاسخ{
	
		 داریم:
	$$\frac{1}{2} \leq \sum_{j = 2^i}^{2^{i+1} - 1} \frac{1}{j} \leq 1$$
	$$\Rightarrow \frac{\floor{\log_2(n)}}{2} \leq T(n) = \sum_{i = 1}^{n} \frac{1}{i} \leq \ceil{\log_2(n)}$$
	$$\Rightarrow T \in \Theta(\log(n))$$
	
}
\زیرمسئله $T(n) = 2T(\sqrt[3]{n}) + \log(n)^2$

\پاسخ{
		از تغییر متغیر $m = log(n)$ استفاده می‌کنیم و سپس قرار می‌دهیم $S(m) = T(2^m)$. در این صورت خواهیم داشت:
	$$S(m) = 2S(\frac{m}{3}) + m^2$$
	طبق قضیه‌ی اصلی $S \in \Theta(m^2)$. پس $T \in \Theta(log(m)^2)$.
	
}

\زیرمسئله \متن‌سیاه{قضیه‌ی اصلی:} در این قسمت مشخص کنید که هر کدام از توابع بازگشتی با قضیه‌ی اصلی قابل تحلیل است یا خیر. در صورتی که تابعی با قضیه‌ی اصلی قابل تحلیل نباشد نیاز به تحلیل آن در این قسمت نیست.
\شروع{فقرات}
\فقره $T(n) = 3T(\frac{n}{5}) + n^{0.6}$
\فقره $T(n) = 3T(\frac{n}{3}) + \frac{n}{log(n)}$
\پایان{فقرات}

\پاسخ{
	
		\شروع{فقرات}
	\فقره $n^{\log_b a} = n^{0.682}$
	و چون 
	$n^{0.6}$
	از  
	$n^{0.682}$
	به صورت چند جمله ای  کوچکتر است پس در این حالت جواب 
	$\theta(n^{0.682})$
	است. (حالت اول قضیه اصلی)
	\فقره $n^{\log_b a} = n^{1}$
	و چون 
	$\frac{n}{log(n)}$
	از  
	$n$
	به صورت چند جمله ای کوچکتر نیست پس در این حالت نمی توان  از قضیه اصلی استفاده کرد. 
	\پایان{فقرات}
	
}
\پایان{زیرمسائل}
%اگر مسئله چند قسمت داشت به این صورت عمل کنید:
%\شروع{زیرمسائل}
%\زیرمسئله[0] % مسئلۀ ۱آ
%\زیرمسئله % مسئلۀ ۱ب
%\زیرمسئله % مسئلۀ ۱پ
%\زیرمسئله % مسئلۀ ۱ت
%\پایان{زیرمسائل}














%%%%%%%%%%%%%%%%%%%%%%%%%%%%%%%%%%%%%%%%%%%%%%%%%%%%%
% برخی دستورات مفید برای تمرین‌ها. %
%%%%%%%%%%%%%%%%%%%%%%%%%%%%%%%%%%%%%%%%%%%%%%%%%%%%%

% فرمول‌های نمونه:
%$f(x) = \Theta(x)$
%$T(x, y) \leq \log(x) + 2^y + \binom{2n}{n}$
% {\tt code\_function}


% لیست موارد بدون شماره‌گذاری:
%\شروع{فقرات}
%    \فقره مورد اول 
%        \شروع{فقرات}
%            \فقره مورد اول داخلی 
%                \شروع{فقرات}
%                    \فقره اولین مورد خیلی داخلی 
%                    \فقره دومین مورد خیلی داخلی 
%                \پایان{فقرات}
%            \فقره مورد دوم داخلی 
%        \پایان{فقرات}
%    \فقره مورد دوم 
%\پایان{فقرات}

% لیست موارد شماره‌گذاری شده:
%\شروع{شمارش}
%    \فقره یک مورد 
%    \فقره یک مورد دیگر 
%    \فقره باز هم یک مورد دیگر
%\پایان{شمارش}

% برای نوشتن یک فرمول چند خطی که یک شماره بخوره. برای حالت‌های دیگر به سایت زیر مراجعه کنید:
% https://www.overleaf.com/learn/latex/Aligning_equations_with_amsmath
%\begin{align} 
%    \begin{split}
%        (x+y)^3 &= (x+y)^2(x+y) \\
%                &= (x^2+2xy+y^2)(x+y) \\
%                &= (x^3+2x^2y+xy^2) + (x^2y+2xy^2+y^3) \\
%                &= x^3+3x^2y+3xy^2+y^3
%    \end{split}                                 
%\end{align}

% نوشتن ماتریس و آرایه دوبعدی
%\begin{align}
%    A = 
%    \begin{bmatrix}
%        A_{11} & A_{21} \\
%        A_{21} & A_{22}
%    \end{bmatrix}
%\end{align}

% اضافه کردن تصویر با شرح. برای اضافه کردن تصویر بدون شرح، فقط دستور درج‌تصویر را استفاده کنید.
%\شروع{شکل}[h!]
%\تنظیم‌ازوسط
%  \درج‌تصویر[پهنا=0.5\پهنای‌متن]{img.jpg}
%\شرح{شرح تصویر}
%\برچسب{تصویر:۱}
%\پایان{شکل}
%\پایان{زیرمسائل}

%کد زیر را ببینید.
%%\lstinputlisting{source_filename.py}
%\begin{latin}
%\begin{lstlisting}
%# Import the modules
%import sys
%import random
%
%ans = True
%
%while ans:
%question = raw_input("Ask the magic 8 ball a question: (press enter to quit) ")
%
%answers = random.randint(1,8)
%
%if question == "":
%sys.exit()
%
%elif answers == 1:
%print "It is certain"
%
%elif answers == 2:
%print "Outlook good"
%
%elif answers == 3:
%print "You may rely on it"
%
%elif answers == 4:
%print "Ask again later"
%
%elif answers == 5:
%print "Concentrate and ask again"
%
%elif answers == 6:
%print "Reply hazy, try again"
%
%elif answers == 7:
%print "My reply is no"
%
%elif answers == 8:
%print "My sources say no"
%\end{lstlisting}
%\end{latin}
